\section{引言}

\subsection{项目背景}

随着金融市场的快速发展，金融研报的数量和复杂度不断增加。分析师和投资者需要从大量的研报中快速提取关键信息，做出明智的投资决策。然而，传统的人工阅读方式效率低下，难以应对海量的信息。

近年来，大语言模型（LLM）和检索增强生成（RAG）技术的发展为金融研报分析提供了新的解决方案。通过结合向量检索和大语言模型的能力，可以实现对金融研报的智能分析和问答。

\subsection{项目目标}

Insight: 金融研报分析引擎项目的核心目标是构建一个对标 NotebookLM 体验的垂直领域 RAG 问答系统，专注于长篇幅、高密度的金融 PDF 研报分析，提供精确到页码的引用溯源。

具体目标包括：

1. 实现金融 PDF 研报的高效解析和处理
2. 构建高质量的向量索引，支持快速准确的信息检索
3. 集成本地大语言模型，提供智能问答和分析能力
4. 实现精确到页码的引用溯源，增强回答的可信度
5. 建立完善的评估体系，确保系统性能

\subsection{技术栈}

本项目采用以下技术栈：

- **前端框架**：Gradio
- **后端框架**：LlamaIndex v0.10+
- **向量数据库**：ChromaDB
- **嵌入模型**：BAAI/bge-m3
- **大语言模型**：Qwen2.5-7B/72B（通过 Ollama 本地部署）
- **PDF 解析**：PyMuPDF
- **评估框架**：RAGAS

\subsection{项目结构}

项目采用模块化设计，主要包含以下目录结构：

- `src/ingest/`：数据接入层，负责 PDF 解析和分块
- `src/retrieval/`：检索层，实现向量检索和混合检索
- `src/generation/`：生成层，负责构建对话引擎和生成回答
- `src/evaluation/`：评估层，实现 RAGAS 评估
- `src/ui/`：交互层，构建 Gradio 前端
- `src/utils/`：工具模块，提供配置管理等功能
- `configs/`：配置文件目录
- `data/`：数据文件目录
- `docs/`：文档目录
